% ============================================================================
\section{Update Sparsity and Teacher Overlap}
\label{sec:subnetworks}
% ============================================================================

We ask whether OPD parameter updates concentrate on a small fraction of parameters,
whether different teachers affect the same parameters within one MOPD
student, and whether teacher distribution differences predict this overlap.
Each question includes its planned experiment; no outcome is assumed.

\subsection{Does update sparsity persist under multiple teachers?}

Prior work reports concentrated parameter changes in RL fine-tuning
\citep{mukherjee2025subnetworks} and OPD
\citep{yu2026denseupdates,shen2026opdgeometry}. We measure cumulative
parameter changes $\Delta=\theta-\theta_0$ and actual optimizer steps
$u=\theta_{\mathrm{after}}-\theta_{\mathrm{before}}$ separately.
For each, we report the fraction of parameters with absolute change at
most a threshold $\tau$, together with the fraction of squared norm carried
by the largest-magnitude $\rho$ fraction of parameters. We sweep $\tau$
and $\rho$: ``update sparsity'' here means near-zero parameter changes or
concentrated changes, not necessarily exact zeros. Per-step and cumulative
sparsity can differ because successive updates can cancel.

\paragraph{Experiment.}
Using the representative single-teacher and joint runs from
Section~\ref{sec:density}, record cumulative changes at early, intermediate,
and final checkpoints and actual optimizer steps around those checkpoints
and initialization. Report curves against both domain exposure and compute.
Measure differences using FP32 master weights when available and report
exported BF16 checkpoint differences separately to identify rounding-induced
zeros; converting BF16 values to FP32 cannot recover lost changes.
Threshold curves, norms, and energy concentration jointly show whether
apparent sparsity reflects small overall scale or concentration on particular
parameters.

As a secondary diagnostic, compare each domain's raw gradient
$g_d=\nabla_\theta\mathcal L_d$ and the combined MOPD gradient under the
same loss reduction, where $\mathcal L_d$ is that domain's mean loss.
Match diagnostic response counts per domain across single-teacher and joint
students, report valid-token
counts, and check sensitivity to batch size; record the combined gradient's
full batch size separately. Repeat over fresh student-rollout batches,
record accumulation denominators, and
accumulate gradients in FP32 before clipping or optimizer processing.
Apply the same threshold, norm, and concentration summaries to gradients,
keeping their scales separate from steps and cumulative changes. AdamW
momentum and weight decay can move parameters with a zero current gradient,
so gradient sparsity alone cannot establish update sparsity.

\subsection{Do teachers affect the same parameters in one student?}

At a fixed MOPD checkpoint $\theta$, compute a proposed optimizer step for
each routed domain using its teacher $\pi_{\phi_d}$ and student-generated
responses. Each calculation starts from an identical copy of the student
and optimizer state, without advancing training. These steps measure the
local effect of applying each teacher's loss; they do not identify fixed,
isolated subnetworks.
To compare locations fairly, first select the same fraction of parameters
for every teacher, ranked by absolute proposed change. For domains $d$ and $e$,
\begin{equation}
\text{overlap}(d,e)=
\frac{\text{number of parameters selected by both teachers}}
     {\text{number of parameters selected by either teacher}}.
\label{subnet_overlap}
\end{equation}
Zero means no shared selected parameters; one means identical selections.
The corresponding difference is one minus this overlap. We also report
absolute-threshold selections, which retain information about update
scale as well as location.

\paragraph{Experiment.}
Use the matched diagnostic response counts and loss reduction above.
Repeat over diagnostic batches and report overlap matrices across
checkpoints, with concentration curves alongside them. Compare global
selections and selections made separately within each layer against
random selections with matching counts, including per-layer results.
Signed step alignment provides additional context: shared parameters
can receive compatible or opposing updates, so lower overlap is not
assumed better.

Separately compare raw-gradient overlap and alignment and use the
zero-gradient optimizer control in Appendix~\ref{app:records} to check
whether shared momentum or decay dominates proposed-step overlap. These
steps do not decompose the joint AdamW step or past changes by teacher.
Independently trained single-teacher endpoint changes provide a reference only: those students
follow different trajectories and cannot identify teacher contributions
inside the joint student. Appendix~\ref{app:records} specifies recording
and threshold conventions.

\subsection{Do more different teachers affect more different parameters?}

We measure the difference between teachers $\pi_{\phi_d}$ and
$\pi_{\phi_e}$ by their Jensen--Shannon divergence, a symmetric, bounded
measure of how differently they distribute probability over next tokens.
Every teacher scores the same cached, held-out student prefixes, with
common domain weights. We also record each teacher's difference from the
current student $\pi_\theta$ and its gradient and step norms.
Appendix~\ref{app:teacher_distance} gives the distribution-distance formula.

\paragraph{Experiment.}
At each checkpoint, plot teacher distribution difference against the
parameter-selection difference from the routed proposed steps above. A smaller
control recomputes teacher steps and gradients on identical prefixes, testing
whether an association persists when the response contexts also match.
For PG, pair the same freshly sampled next tokens and fixed prefix weights
across teachers, repeating the paired draws across batches.
Available intermediate teacher checkpoints can add variation without
training new teachers.

Show individual pair trajectories and run-level uncertainty. Pairs that
share a teacher or training lineage, and repeated checkpoints, are
dependent observations. Relate teacher--student differences to update
concentration and observed transfer as additional analyses. A small
teacher pool supports a setting-specific association, not a universal
monotone relationship; larger, smaller, or unchanged parameter overlap
are all valid outcomes.
